\documentclass[11pt]{article}

\usepackage[margin=1in]{geometry}
\usepackage{amsmath,amssymb}
\usepackage{booktabs}
\usepackage{longtable}
\usepackage{array}
\usepackage{tabularx}
\usepackage{multirow}
\usepackage{siunitx}
\usepackage{xcolor}
\usepackage[hidelinks]{hyperref}
\usepackage{fancyhdr}

\title{Reference Benchmarks for the CoQui PAW + ISDF--THC\\
       Validation Suite: Band Gaps and RPA Energies}
\author{Compiled from the published literature}
\date{April 30, 2026}

\pagestyle{fancy}
\fancyhf{}
\fancyhead[L]{CoQui PAW Validation References}
\fancyhead[R]{\thepage}

\newcommand{\eV}{\ensuremath{\,\mathrm{eV}}}
\newcommand{\meV}{\ensuremath{\,\mathrm{meV}}}
\newcommand{\GPa}{\ensuremath{\,\mathrm{GPa}}}
\newcommand{\target}[1]{\textbf{\textcolor{blue!60!black}{#1}}}

\begin{document}

\maketitle

\begin{abstract}
\noindent
This document is a reference compendium for the CoQui PAW + ISDF--THC
validation suite. It collects published quasiparticle band gaps (G$_0$W$_0$,
GW$_0$, scGW) and RPA total/correlation energies for benchmark solids,
spanning VASP, ABINIT, BerkeleyGW, GPAW, FHI--aims, exciting, SPEX/FLEUR
and Quantum~ESPRESSO. Two systems anchor the validation: Si (the standard
``easy'' benchmark) and ZnO (the well-known ``hard'' benchmark). Secondary
systems --- diamond, MgO, LiF, NaCl, GaAs, hBN, graphite, noble-gas solids,
and selected metals --- are included for completeness. For each entry we
record the level of theory, code, projector type (PAW / ONCV / all-electron
LAPW / NAO), k-mesh, band cutoff, frequency treatment, and the source DOI.
The document also flags published convergence pitfalls (notably the
ZnO ``false-convergence trap'' and the basis-set $E_{\mathrm{cut}}^{-3/2}$
tail in RPA) so they can be avoided when comparing CoQui results to the
literature.
\end{abstract}

\tableofcontents

\newpage

\section{Executive summary}

\subsection{What to target for CoQui}

\begin{table}[h]
\centering
\caption{Headline targets for the first round of CoQui PAW+ISDF--THC
validation. Numbers are typical converged values from the literature; see
the body of the document for full tables and sources.}
\label{tab:targets}
\small
\begin{tabular}{l l c c c}
\toprule
System & Quantity                 & Target value                   & Reference         & Difficulty \\
\midrule
Si     & G$_0$W$_0$@PBE indirect gap & $\approx 1.1$--$1.2\eV$       & multi-code 2024   & easy \\
Si     & RPA@PBE $E_{\mathrm{coh}}$  & $\approx 4.6\eV$/atom          & Harl 2010 (VASP)  & easy \\
Si     & RPA@PBE lattice constant   & $5.421\AA$                     & Harl 2010 (VASP)  & easy \\
ZnO    & G$_0$W$_0$@PBE direct gap   & $2.5$--$2.9\eV$                & multi-code 2024   & hard \\
ZnO    & RPA@PBE lattice $a$        & $\approx 3.254\AA$             & Harl 2010 (VASP)  & hard \\
\bottomrule
\end{tabular}
\end{table}

\subsection{Three things to watch}

\begin{enumerate}
  \item \textbf{Always extrapolate $E_{\mathrm{cut}}^{-3/2}$ for RPA.}
  Single-cutoff RPA correlation energies are not comparable across codes;
  the tail is large (0.3--0.5\,eV/atom for Si at typical cutoffs).
  \item \textbf{ZnO needs Zn 3s/3p/3d in valence and large virtual space}
  for both GW and RPA. Anything less than 1000 bands plus extrapolation
  is the well-documented false-convergence trap.
  \item \textbf{Mind the starting point.} RPA@LDA and RPA@PBE differ by
  $\sim 0.2$--$0.3\eV$/atom for covalent solids; G$_0$W$_0$@LDA and
  G$_0$W$_0$@PBE/PBEsol differ by $\sim 0.1\eV$ for Si. Compare
  apples-to-apples.
\end{enumerate}

\subsection{Bottom-line accuracy expectations}

For Si, the four-code 2024 benchmark (Borlido et al., arXiv:2411.19701)
gives a 112\,meV cross-code spread for the G$_0$W$_0$@PBEsol gap. For ZnO
the same benchmark shows 334\,meV spread --- the largest of any solid in
the test set. Anything CoQui produces inside the 2.5--2.9\,eV window for
ZnO with proper convergence is consistent with the established
literature; values below $\sim 2.0\eV$ indicate a PAW augmentation or
basis-set issue.

\newpage
\section{Part I --- GW band gaps}

\subsection{Silicon (Si, diamond structure)}

\textbf{Structure:} Fd--3m, $a = 5.431\AA$. The gap is indirect
$\Gamma \to \Delta$ (about $0.85 \times X$). Experimental indirect gap:
\target{1.17\eV} (0\,K, no zero-point renormalization).

\begin{longtable}{p{3.0cm} c l l l l p{3.5cm}}
\caption{G$_0$W$_0$, GW$_0$, scGW gap of Si.} \label{tab:si_gw} \\
\toprule
Level & Gap (eV) & Code & PP & k-mesh & Freq. & Source \\
\midrule
\endfirsthead
\multicolumn{7}{l}{\emph{Table~\ref{tab:si_gw} continued.}} \\
\toprule
Level & Gap (eV) & Code & PP & k-mesh & Freq. & Source \\
\midrule
\endhead
LDA (KS)         & 0.51  & VASP    & PAW       & ---     & ---       & textbook \\
G$_0$W$_0$@LDA   & 0.85  & VASP    & PAW       & 4$\times$4$\times$4 & PP (GN) & Shishkin/Kresse 2007 \\
GW$_0$@LDA       & 1.22  & VASP    & PAW       & 4$\times$4$\times$4 & PP    & Shishkin/Kresse 2007 \\
scGW@LDA         & 1.41  & VASP    & PAW       & 4$\times$4$\times$4 & PP    & Shishkin/Kresse 2007 \\
G$_0$W$_0$@LDA   & 1.094 & VASP    & PAW       & 10$\times$10$\times$10 & CD & VASP wiki tutorial \\
G$_0$W$_0$@PBEsol& 1.257 & ABINIT  & ONCV      & 8$\times$8$\times$8 & CD+AC & Borlido 2024 \\
G$_0$W$_0$@PBEsol& 1.145 & exciting& LAPW+lo   & 8$\times$8$\times$8 & 32 imag.\ freq. & Borlido 2024 \\
G$_0$W$_0$@PBEsol& 1.147 & FHI-aims& NAO (all-e)& 8$\times$8$\times$8 & 60 freq. (Pad\'e) & Borlido 2024 \\
G$_0$W$_0$@PBEsol& 1.152 & GPAW    & PW (NC-PP)& 8$\times$8$\times$8 & real-freq.\ extrap.& Borlido 2024 \\
\target{Experiment} & \target{1.17} & --- & --- & --- & --- & standard \\
\bottomrule
\end{longtable}

\noindent \emph{Highlights.}
G$_0$W$_0$@LDA underestimates by $\sim 0.3\eV$. Eigenvalue self-consistency
in $G$ (GW$_0$) closes the gap to within $0.05\eV$ of experiment. The
2024 multi-code benchmark with PBEsol orbitals gives 1.145--1.257\,eV;
all-electron codes (FHI-aims, exciting) cluster at $\sim 1.15\eV$. The
ABINIT outlier at 1.257\,eV reflects residual ONCV / $\epsilon_{\mathrm{cut}}$
sensitivity.

\subsection{ZnO (wurtzite)}

\textbf{Structure:} P6$_3$mc, $a = 3.250\AA$, $c = 5.207\AA$.
The gap is \emph{direct} $\Gamma \to \Gamma$. Experimental optical gap:
\target{3.44\eV} (4\,K). Zn 3d band at $\sim -7\eV$. LDA gives
$\sim 0.7$--$0.8\eV$.

\begin{longtable}{p{3.2cm} c l l c c l p{3.0cm}}
\caption{G$_0$W$_0$ gap of ZnO --- the canonical hard benchmark.}
\label{tab:zno_gw} \\
\toprule
Level & Gap (eV) & Code & PP & Zn 3d & k-mesh & Freq. & Source \\
\midrule
\endfirsthead
\multicolumn{8}{l}{\emph{Table~\ref{tab:zno_gw} continued.}} \\
\toprule
Level & Gap (eV) & Code & PP & Zn 3d & k-mesh & Freq. & Source \\
\midrule
\endhead
LDA (KS)              & 0.76     & ABINIT    & ONCV  & val.\ & 8$\times$8$\times$5 & ---       & Borlido 2024 \\
G$_0$W$_0$@LDA        & $\sim$0.8 & VASP     & PAW   & core  & 4$\times$4$\times$4 & PP & Shishkin/Kresse 2007 \\
GW$_0$@LDA            & $\sim$2.97 & VASP    & PAW   & core  & 4$\times$4$\times$4 & PP & Shishkin/Kresse 2007 \\
G$_0$W$_0$@LDA conv.\ & $\sim$2.5 & BerkeleyGW & NC-PP & val.\ & 4$\times$4$\times$3 & GPP & Shih 2010 \\
G$_0$W$_0$@LDA conv.\ & 3.0--3.2 & SPEX     & FLAPW & val.\ & 4$\times$4$\times$4 & full-freq.\ & Friedrich 2011 \\
G$_0$W$_0$@LDA (CD)   & $\sim$2.4 & ABINIT   & PAW   & val.\ & 6$\times$6$\times$4 & CD & Stankovski 2011 \\
G$_0$W$_0$@LDA+U      & $\sim$3.2 & VASP     & PAW   & val.\ & 6$\times$6$\times$4 & PP & Lany/Zunger 2010 \\
G$_0$W$_0$@PBEsol     & 2.613    & ABINIT    & ONCV  & val.\ & 8$\times$8$\times$5 & CD+AC & Borlido 2024 \\
G$_0$W$_0$@PBEsol     & 2.878    & exciting  & LAPW  & val.\ & 8$\times$8$\times$5 & 32 imag. & Borlido 2024 \\
G$_0$W$_0$@PBEsol     & 2.758    & FHI-aims  & NAO   & val.\ & 8$\times$8$\times$5 & 60 freq. & Borlido 2024 \\
G$_0$W$_0$@PBEsol     & 2.544    & GPAW      & PW    & val.\ & 8$\times$8$\times$5 & real-freq.\ & Borlido 2024 \\
\target{Experiment (optical)} & \target{3.44} & --- & --- & --- & --- & --- & standard \\
\target{Experiment (transport)} & \target{3.37} & --- & --- & --- & --- & --- & standard \\
\bottomrule
\end{longtable}

\noindent \emph{The convergence saga.}
The published ZnO G$_0$W$_0$ value drifted from $\sim 0.8\eV$ (2007) to
$\sim 2.5$--$2.9\eV$ (2024) as two issues were resolved. (1) Zn 3s/3p/3d
must be in valence, not core. (2) The number of empty bands must be
$\gg 1000$ with explicit $E_{\mathrm{cut}}^{-3/2}$ extrapolation, because
the Coulomb-hole self-energy $\Sigma_{\mathrm{COH}}$ shows
\emph{non-uniform convergence}: VBM and CBM converge at different rates,
biasing the gap. f-sum-rule plasmon-pole models are unreliable for ZnO
because the shallow Zn 3d resonance pollutes the f-sum; contour
deformation or full-frequency integration is required.

\subsection{Secondary GW benchmarks}

\paragraph{MgO (rock-salt, Fm--3m).}
Direct $\Gamma \to \Gamma$ gap; experiment \target{7.83\eV}; LDA $\sim 4.5\eV$.
VASP PAW (Shishkin/Kresse 2007): G$_0$W$_0$@LDA $\sim 6.2$--$6.5\eV$,
GW$_0$@LDA $\sim 7.2$--$7.5\eV$, scGW@LDA $\sim 7.8\eV$.

\paragraph{LiF (rock-salt, Fm--3m).}
Direct gap; experiment \target{14.2\eV}; LDA $\sim 9.2\eV$.
VASP PAW: G$_0$W$_0$@LDA 12.0--12.5\,eV, GW$_0$@LDA $\sim 13.5\eV$,
scGW@LDA $\sim 14.8\eV$ (overshoots).

\paragraph{GaAs (zinc-blende, F--43m).}
Direct gap, experiment \target{1.52\eV} at 0\,K; LDA $\sim 0.2$--$0.5\eV$;
Ga 3d at $\sim -18\eV$. VASP PAW (Ga 3d as core): G$_0$W$_0$@LDA
$\sim 1.1$--$1.3\eV$, GW$_0$@LDA $\sim 1.5\eV$. BerkeleyGW (Ga 3d valence):
G$_0$W$_0$@LDA 1.3--1.5\,eV.

\paragraph{TiO$_2$ rutile (P4$_2$/mnm).}
Direct optical gap, experiment \target{3.03\eV}; LDA $\sim 1.8\eV$.
2024 cross-code G$_0$W$_0$@PBEsol: ABINIT 3.333, exciting 3.321,
FHI-aims 3.546, GPAW 3.281\,eV.

\subsection{Cross-code precision benchmark (2024)}

\textbf{Source:} Borlido et~al., \emph{arXiv:2411.19701} /
Comput.~Mater.~Sci.\ 2025, DOI 10.1016/j.commatsci.2024.113495.
PBEsol orbitals; full-frequency methods (no plasmon-pole). Zn valence in
ABINIT ONCV is 20 electrons; exciting and FHI-aims are all-electron.

\begin{table}[h]
\centering
\caption{Four-code G$_0$W$_0$@PBEsol gaps for Si, ZnO, TiO$_2$
(eV). Spread is the max--min range across the four codes.}
\label{tab:multicode_2024}
\small
\begin{tabular}{l c c c c c c}
\toprule
Material & ABINIT & exciting & FHI-aims & GPAW & Spread (meV) & Expt \\
\midrule
Si (indirect)        & 1.257 & 1.145 & 1.147 & 1.152 & 112 & 1.17 \\
ZnO (direct)         & 2.613 & 2.878 & 2.758 & 2.544 & 334 & 3.44 \\
TiO$_2$ rutile (dir.) & 3.333 & 3.321 & 3.546 & 3.281 & 265 & 3.03 \\
\bottomrule
\end{tabular}
\end{table}

\subsection{Six-step ZnO convergence checklist}

\begin{enumerate}
  \item \textbf{Semicore in valence.} Zn 3s, 3p, 3d in the active space
        (12 vs 20 valence electrons). Core treatment $\Rightarrow$
        $\sim 0.8\eV$ gap, wrong by $\sim 2.6\eV$.
  \item \textbf{Empty bands.} $\sim 1000$--$3000$ bands or basis-set
        extrapolation; $\Sigma_{\mathrm{COH}}$ converges non-uniformly.
  \item \textbf{$\boldsymbol{\chi}$ cutoff.} ENCUTGW $\gtrsim 30$--$50$\,Ha
        in plane-wave codes; otherwise a false convergence plateau.
  \item \textbf{Frequency scheme.} Avoid f-sum-rule plasmon-pole;
        use contour deformation or full-frequency integration.
  \item \textbf{Starting point.} LDA places Zn 3d too high
        ($-5$ vs $-7\eV$). LDA+U or hybrid starting points improve
        convergence.
  \item \textbf{PAW partial-wave completeness.} Standard PAW partial waves
        are incomplete for high-energy unoccupied states; norm-conserving
        partial waves or basis-set corrections (Klime\v{s}/Kaltak/Kresse,
        PRB 90, 075125, 2014) are needed for predictive results.
\end{enumerate}

\newpage
\section{Part II --- RPA total / correlation energies}

\textbf{Convention.} In what follows ``RPA'' means RPA correlation on top
of a DFT exchange piece, i.e.\ $E_{\mathrm{tot}} = E_x[\phi^{\mathrm{KS}}]
+ E_c^{\mathrm{RPA}}[\chi_0]$. ``RPA@PBE'' is the standard choice in the
literature (Harl/Kresse). Cohesive energies are reported as positive
$=$ bound, in eV/atom unless noted.

\subsection{Silicon}

Experimental cohesive energy: \target{4.62\eV/atom} (0\,K, ZPE-corrected).
Lattice constant: \target{5.431\AA}; bulk modulus \target{99\GPa}.

\begin{table}[h]
\centering
\caption{Si cohesive energy at the RPA level.}
\label{tab:si_rpa_ecoh}
\small
\begin{tabular}{l c l l l p{4cm}}
\toprule
Method & $E_{\mathrm{coh}}$ (eV/atom) & Code & PP & k-mesh & Source \\
\midrule
PBE                  & 4.55     & VASP  & PAW & 8$\times$8$\times$8 & Harl 2010 \\
RPA@PBE              & $\sim$4.6 & VASP & PAW & 8$\times$8$\times$8 & Harl 2010 \\
RPA@LDA              & 4.32     & GPAW  & PAW & 4$\times$4$\times$4 & Olsen/Thygesen 2013 \\
EXX@PBE (no corr.)   & 2.82     & GPAW  & PAW & 4$\times$4$\times$4 & GPAW tutorial \\
\target{Experiment}  & \target{4.62} & --- & --- & --- & standard \\
\bottomrule
\end{tabular}
\end{table}

\begin{table}[h]
\centering
\caption{Si lattice constant and bulk modulus, RPA vs PBE vs experiment.}
\label{tab:si_rpa_struct}
\small
\begin{tabular}{l c c l l}
\toprule
Method & $a$ (\AA) & $B$ (GPa) & Code & Source \\
\midrule
PBE              & 5.468 & 90    & VASP/PAW   & Harl 2010 \\
RPA@PBE          & 5.421 & $\sim$98 & VASP/PAW & Harl 2010 \\
RPA              & 5.414 & ---   & QE/ONCV    & Pitts 2025 \\
PBE              & 5.479 & ---   & QE/ONCV    & Pitts 2025 \\
\target{Expt.}   & \target{5.431} & \target{99} & --- & standard \\
\bottomrule
\end{tabular}
\end{table}

\subsection{Diamond, MgO, LiF, NaCl}

Compact summary; full per-system tables in
\texttt{notes/paw\_rpa\_reference\_benchmarks.md}.

\begin{table}[h]
\centering
\caption{RPA@PBE structural and cohesive properties (VASP/PAW, Harl 2010,
unless noted). Experimental values in the right column.}
\label{tab:secondary_rpa}
\small
\begin{tabular}{l c c c c c c}
\toprule
Material   & $a_{\mathrm{PBE}}$ & $a_{\mathrm{RPA}}$ & $a_{\mathrm{expt}}$ & $E_{\mathrm{coh}}^{\mathrm{RPA}@\mathrm{PBE}}$ & $E_{\mathrm{coh}}^{\mathrm{expt}}$ & Notes \\
           & (\AA)  & (\AA)  & (\AA)  & (eV/at.) & (eV/at.) & \\
\midrule
Si              & 5.468 & 5.421 & \target{5.431} & $\sim 4.6$    & \target{4.62} & covalent \\
C (diamond)     & 3.572 & 3.553 & \target{3.553} & $\sim 7.4$--$7.5$ & \target{7.37} & covalent, $B=443\GPa$ \\
BN (zinc-blende)& 3.626 & 3.604 & \target{3.592} & ---           & ---           & --- \\
MgO             & 4.258 & 4.196 & \target{4.211} & $\sim 5.0$--$5.1$ & \target{5.20} & ionic \\
LiF             & 4.073 & 3.986 & \target{3.972} & $\sim 4.3$--$4.4$ & \target{4.46} & ionic \\
NaCl            & 5.694 & 5.573 & \target{5.595} & $\sim 3.3$    & \target{3.34} & soft ionic \\
ZnO ($a$)       & 3.283 & 3.254 & \target{3.250} & ---           & $\sim 3.60$/atom & semicore \\
ZnO ($c$)       & 5.304 & 5.262 & \target{5.207} & ---           & ---           & semicore \\
\bottomrule
\end{tabular}
\end{table}

\subsection{Layered and weakly-bound systems}

\paragraph{Graphite interlayer binding.}
The landmark RPA-for-graphite paper (Leb\`egue et~al., PRL 105, 196401,
2010): VASP/PAW RPA@PBE gives an interlayer binding energy of
\target{48\,meV/atom} at $d_{\mathrm{int}} = 3.34\AA$, in agreement with
the experimental window 35--52\,meV/atom. PBE gives essentially zero
interlayer binding. RPA captures the correct $1/d^3$ van~der~Waals
asymptotics.

\paragraph{Hexagonal BN.}
RPA@PBE in-plane $a \approx 2.50\AA$, interlayer binding 30--40\,meV/atom
(VASP/PAW, Harl 2010). PBE: no binding.

\paragraph{Noble-gas solids.}
Harl \& Kresse, PRB 77, 045136 (2008), VASP/PAW, RPA@LDA. The foundational
RPA-for-vdW-solids paper.

\begin{table}[h]
\centering
\caption{Noble-gas solids: lattice constant and cohesive energy at RPA@LDA
vs experiment.}
\label{tab:noble_gas}
\small
\begin{tabular}{l c c c c}
\toprule
System & $a_{\mathrm{RPA}}$ (\AA) & $a_{\mathrm{expt}}$ (\AA) & $E_{\mathrm{coh}}^{\mathrm{RPA}}$ (meV/at) & $E_{\mathrm{coh}}^{\mathrm{expt}}$ (meV/at) \\
\midrule
Ne & $\sim 4.46$ & 4.430 & $\sim 27$  & 27 \\
Ar & $\sim 5.26$ & 5.256 & $\sim 88$  & 88 \\
Kr & $\sim 5.64$ & 5.646 & $\sim 120$ & 116 \\
\bottomrule
\end{tabular}
\end{table}

\subsection{Metals}

Schimka et~al., \emph{Nature Materials} 9, 741 (2010) and PRB 87, 214102
(2013). RPA@PBE cohesive energies for bulk metals are roughly
PBE-quality; the gain from RPA is concentrated in surface and adsorption
energies, not bulk cohesion. Approximate values (eV/atom): Cu $\sim 3.5$,
Fe $\sim 4.3$, Na $\sim 1.1$, Mg $\sim 1.5$, Al $\sim 3.4$.

\subsection{Statistical summary across 20 solids (Harl 2010)}

\begin{itemize}
  \item PBE lattice constant MAE $\approx 0.05\AA$ (systematic
        overestimate);\\
        RPA@PBE lattice constant MAE $\approx 0.02\AA$ (mixed).
  \item PBE bulk modulus MAE $\approx 9\GPa$;
        RPA@PBE MAE $\approx 5\GPa$.
  \item RPA@PBE is a clear improvement over PBE for structural
        properties.
  \item RPA \emph{underbinds} covalent solids (Si, C) by 0.1--0.3\,eV/atom
        relative to experiment; ionic solids fare better.
\end{itemize}

\subsection{Convergence parameters by code}

\begin{table}[h]
\centering
\caption{Recommended RPA convergence parameters by code.}
\label{tab:rpa_conv}
\small
\begin{tabular}{l p{2.4cm} p{2.4cm} p{2.6cm} p{4cm}}
\toprule
Code & Cutoff for $\chi_0$ & N$_\omega$ & N$_\mathrm{bands}$ & Extrapolation \\
\midrule
VASP/PAW   & ENCUTGW = ENCUT (or 2/3) & 12--16 Gauss-Legendre & $\sim 3$--$5\times N_{\mathrm{occ}}$ & $E_{\mathrm{cut}}^{-3/2}$ on $E_{\mathrm{RPA}}$ \\
GPAW/PAW   & 80--164\,eV ($\chi_0$ cutoff) & 16 GL & set by N$_{\mathrm{pw}}$ & $E_{\mathrm{cut}}^{-3/2}$ \\
FHI-aims   & all-electron NAO tier-2/3 & 40--60 imag.\ freq. & all virtual NAO & NAO basis tier \\
QE/ONCV    & $\sim 80\,$Ry        & 8 imag.\ freq. & $\sim 5\times N_e$ & basis ext.\ implicit \\
\bottomrule
\end{tabular}
\end{table}

\paragraph{The $E_{\mathrm{cut}}^{-3/2}$ tail.}
The leading correction to a finite-cutoff RPA correlation energy is
$E_{\mathrm{RPA}}(E_{\mathrm{cut}}) \approx E_{\mathrm{RPA}}(\infty)
+ A\, E_{\mathrm{cut}}^{-3/2}$ (Harl/Kresse 2008). For Si at
ENCUT $\sim 400$\,eV the missing tail is 0.3--0.5\,eV/atom. \emph{Always
extrapolate.} Single-cutoff numbers are not portable across codes.

\newpage
\section{Recommended validation order for CoQui}

\begin{enumerate}
  \item \textbf{Si} --- both G$_0$W$_0$ and RPA. The simplest check on
        the full PAW + ISDF--THC pipeline. Targets: G$_0$W$_0$@PBE gap
        $\approx 1.1\eV$, RPA@PBE $E_{\mathrm{coh}} \approx 4.6\eV$/atom,
        $a_{\mathrm{RPA}} \approx 5.421\AA$.
  \item \textbf{MgO and LiF} --- ionic solids where RPA performs well.
        Tests Hartree/exchange balance in the PAW augmentation. Targets:
        $a_{\mathrm{MgO}} \approx 4.196$--$4.200\AA$,
        $a_{\mathrm{LiF}} \approx 3.986$--$3.996\AA$.
  \item \textbf{Diamond} --- covalent, small cell, well-studied. Target
        $a \approx 3.553\AA$, RPA@PBE $E_{\mathrm{coh}} \approx
        7.4$--$7.5\eV$/atom.
  \item \textbf{ZnO} --- semicore-state stress test. Zn 3s/3p/3d in
        valence, $\geq 1000$ bands with extrapolation, contour
        deformation. Targets: G$_0$W$_0$@PBE in 2.5--2.9\,eV window;
        $a_{\mathrm{RPA}} \approx 3.254\AA$. Anything below
        $\sim 2.0\eV$ for the ZnO gap indicates a PAW augmentation or
        basis-set issue.
\end{enumerate}

\section{Key references}

\begin{enumerate}
  \item Shishkin \& Kresse, \emph{Phys.\ Rev.\ B} \textbf{75}, 235102
        (2007). DOI 10.1103/PhysRevB.75.235102. Canonical VASP PAW
        G$_0$W$_0$/GW$_0$/scGW for Si, ZnO, MgO, LiF, GaAs, GaN.

  \item Shih, Xue, Zhang, Cohen, Louie, \emph{Phys.\ Rev.\ Lett.}
        \textbf{105}, 146401 (2010). DOI
        10.1103/PhysRevLett.105.146401. BerkeleyGW; resolves the ZnO
        gap underestimate (semicore + 3000 bands).

  \item Friedrich, M\"{u}ller, Bl\"{u}gel, \emph{Phys.\ Rev.\ B}
        \textbf{83}, 081101(R) (2011); erratum \textbf{84}, 039906.
        DOI 10.1103/PhysRevB.83.081101. SPEX/FLAPW; ZnO as ``extreme
        case''; linearization-error correction; converged $\sim 2.83\eV$.

  \item Stankovski et~al., \emph{Phys.\ Rev.\ B} \textbf{84}, 241201(R)
        (2011). DOI 10.1103/PhysRevB.84.241201. ABINIT; PP models
        unreliable for ZnO; CD gives $\sim 2.4\eV$.

  \item Klime\v{s}, Kaltak, Kresse, \emph{Phys.\ Rev.\ B} \textbf{90},
        075125 (2014). DOI 10.1103/PhysRevB.90.075125. PAW partial-wave
        completeness; NC partial waves + finite-basis correction.

  \item Borlido et~al., arXiv:2411.19701 / \emph{Comput.\ Mater.\ Sci.}
        2025. DOI 10.1016/j.commatsci.2024.113495. Four-code G$_0$W$_0$
        precision benchmark (ABINIT, exciting, FHI-aims, GPAW); Si, ZnO,
        TiO$_2$. \textbf{Most current multi-code reference.}

  \item Harl \& Kresse, \emph{Phys.\ Rev.\ B} \textbf{77}, 045136 (2008).
        DOI 10.1103/PhysRevB.77.045136. Foundational RPA-for-solids
        paper: VASP/PAW ACFDT for Ne, Ar, Kr; establishes
        $E_{\mathrm{cut}}^{-3/2}$ basis extrapolation.

  \item Harl, Schimka, Kresse, \emph{Phys.\ Rev.\ B} \textbf{81},
        115126 (2010). DOI 10.1103/PhysRevB.81.115126.
        \textbf{Primary RPA benchmark for solids:} VASP/PAW RPA@PBE,
        20 solids, lattice constants, bulk moduli, atomization energies.

  \item Schimka et~al., \emph{Nat.\ Mater.} \textbf{9}, 741 (2010).
        DOI 10.1038/nmat2806. RPA@PBE for metals + insulators; surface
        and adsorption energies.

  \item Leb\`egue et~al., \emph{Phys.\ Rev.\ Lett.} \textbf{105},
        196401 (2010). DOI 10.1103/PhysRevLett.105.196401. VASP/PAW RPA
        for graphite interlayer binding (48\,meV/atom).

  \item Olsen \& Thygesen, \emph{Phys.\ Rev.\ B} \textbf{86}, 081103(R)
        (2012); \textbf{87}, 075111 (2013). DOI
        10.1103/PhysRevB.86.081103, 10.1103/PhysRevB.87.075111.
        GPAW RPA benchmarks.

  \item Kaltak, Klime\v{s}, Kresse, \emph{J.\ Chem.\ Theory Comput.}
        \textbf{10}, 2498 (2014). DOI 10.1021/ct5001268. Low-scaling
        $O(N^3)$ RPA via imaginary time + Laplace; minimax quadrature.

  \item Ren, Rinke, Joas, Scheffler, \emph{J.\ Mater.\ Sci.} \textbf{47},
        7447 (2012). DOI 10.1007/s10853-012-6570-4. FHI-aims all-electron
        RPA review with benchmark table for 11 solids.

  \item Booth, Gr\"{u}neis, Kresse, Alavi, \emph{Nature} \textbf{493},
        365 (2013). DOI 10.1038/nature11770. FCIQMC for solids; exact
        wavefunction benchmark vs RPA and CCSD(T).

  \item Pitts, Contant, Hellgren, \emph{Phys.\ Rev.\ B} \textbf{112},
        085137 (2025). DOI 10.1103/PhysRevB.112.085137 / arXiv:2504.12768.
        Self-consistent RPA in QE/ONCV; cross-validation of VASP PAW
        within $0.01\AA$.

  \item Schimka et~al., \emph{Phys.\ Rev.\ B} \textbf{87}, 214102 (2013).
        DOI 10.1103/PhysRevB.87.214102. RPA@PBE lattice constants and
        cohesive energies of alkali / alkaline-earth / transition metals.
\end{enumerate}

\section{Provenance and uncertainty}

The numerical values in this document are drawn from the cited primary
references where the published tables are accessible, and from review
abstracts / cross-citations in newer papers where the original tables
are paywalled (notably Harl 2010 PRB 81, 115126). Values quoted with a
``$\sim$'' carry approximately $\pm 0.1\eV$ (or $\pm 0.01\AA$ for lattice
constants) of indirect-sourcing uncertainty. The DOIs are included so
the original tables can be consulted when institutional access is
available. The companion markdown files
(\texttt{notes/paw\_gw\_reference\_benchmarks.md} and
\texttt{notes/paw\_rpa\_reference\_benchmarks.md}) include longer
annotations and per-row sources.

\end{document}
